\chapter{Unreal Engine 5 a AI systém}

Tato kapitola se zabývá nástroji Unreal Engine 5, které slouží k implementaci chování NPC. 
Nejprve je stručně vysvětlen pojem herního enginu a představen Unreal Engine 5. 
Následně jsou popsány konkrétní systémy, které jsou v této práci využity pro řízení chování postav.

\section{Herní engine a Unreal Engine 5}

Herní engine je softwarový nástroj, který poskytuje základní funkcionalitu potřebnou pro vývoj her. 
Typicky zahrnuje subsystémy pro vykreslování grafiky, fyzikální simulaci, práci se vstupy nebo správu herní logiky. 
Díky tomu se vývojář nemusí zabývat implementací těchto základních částí od začátku a může se soustředit na konkrétní obsah hry.\cite[s.~11--13]{gregory2018} 
Herní engine tedy pomáhá zvýšit potenciální produktivitu vývojáře.

Unreal Engine 5 (UE5) je moderní herní engine vyvíjený společností Epic Games. 
Kromě samotného běhu hry poskytuje také nástroje pro tvorbu herního obsahu a logiky. 
Jednou z jeho výhod je integrace systémů pro tvorbu umělé inteligence, které umožňují navrhovat chování NPC přímo 
v rámci enginu bez nutnosti použití externích knihoven.\cite{epic_gameplay_systems}

V UE5 je možné implementovat herní prvky a logiku buď pomocí C++, nebo pomocí proprietárního vizuálního skriptovacího systému, který se nazývá Blueprint.
V této práci jsou všechny systémy implementovány pomocí Blueprintu.


\section{Pawn a AI Controller}

V Unreal Engine 5 je objekt reprezentující postavu ve hře označován jako Pawn. 
Pawn určuje fyzickou přítomnost postavy ve scéně, například její polohu nebo pohyb.\cite{epic_pawn} 

Každý Pawn má controller, kterým je řízen. 
Pokud je postava ovládána hráčem, používá se Player Controller, v případě NPC pak AI Controller.\cite{epic_controller}

Typicky při převzetí kontroly nad Pawnem spouští behaviorální strom, který definuje jednotlivé akce NPC. 
Standardně platí mezi Pawn a Controller vazba 1:1, ale možná je i vícenásobná vazba (např. pro hry žánru RTS).
Propojení mezi AI Controllerem a Pawnem umožňuje controlleru ovládat pohyb i další aspekty chování postavy.

\section{Behaviorální strom v UE5}

Behaviorální strom je v Unreal Engine implementován jako samostatný asset, který lze upravovat pomocí grafického editoru.\cite{epic_bt} 
Princip jeho fungování odpovídá obecné teorii popsané v předchozí kapitole.

Strom se skládá z uzlů, které definují jednotlivé části chování. 
V UE5 se používají zejména tři typy uzlů:
\begin{itemize}
    \item \textit{Task} – představuje konkrétní akci (např. pohyb na cílovou pozici),
    \item \textit{Service} – pravidelně aktualizuje data během vykonávání větve,
    \item \textit{Decorator} – určuje podmínky, za kterých může být větev vykonána.
\end{itemize}

Behaviorální strom pracuje s daty uloženými v Blackboard a na jejich základě rozhoduje, jaká větev bude v daném okamžiku vykonána. 
Výhodou tohoto přístupu je přehlednost a možnost jednoduché úpravy chování bez zásahu do zdrojového kódu \cite[s.~80]{butler2023}.

\section{Blackboard}

Blackboard slouží jako sdílené datové úložiště pro behaviorální strom. 
Uchovává hodnoty přiřazené ke klíčům, například referenci na hráče, cílovou pozici nebo aktuální stav postavy.\cite{epic_bt} 

Samotný Blackboard neobsahuje žádnou logiku, pouze zpřístupňuje data ostatním částem systému. 
Tyto hodnoty mohou být nastavovány z různých míst, například z behaviorálního stromu nebo přímo z Blueprintu postavy.

\section{Navigační systém}

Pro pohyb NPC v prostředí využívá UE5 navigační systém založený na tzv. Navigation Mesh (NavMesh). 
NavMesh reprezentuje průchodné oblasti scény a určuje, kde se může postava pohybovat.\cite{epic_navsystem} 

Při vykonání akce pohybu (např. pomocí uzlu \textit{MoveTo}) je cílová pozice předána navigačnímu systému, který vypočítá cestu po NavMesh. 
Samotný pohyb postavy je následně realizován komponentou pro pohyb postavy.

\section{AI Perception}

AI Perception umožňuje NPC vnímat okolní prostředí prostřednictvím definovaných smyslů, například zraku nebo sluchu \cite{epic_perception}. 

Při zaznamenání podnětu (např. detekce hráče) dojde k vyvolání události, na jejímž základě se aktualizují hodnoty v Blackboardu. 
Behaviorální strom poté při dalším vyhodnocení reaguje na novou situaci a zvolí odpovídající chování.
